\section{Polynomial-time decision procedures}\label{sec:polytime}

Throughout this section $\Gamma=\langle M_1,\ldots,M_k\rangle\le\GL_d(\Q)$, where
$M_1,\ldots,M_k$ are the minimal letter matrices of an invertible series (computed in
polynomial time by Theorem~\ref{thm:decision}), $\cG$ is the Zariski closure of $\Gamma$, and
$\sigma$ is the total bit length of the entries of the $M_a^{\pm1}$. We work in the Turing model
with exact rational arithmetic. The \emph{height} $\operatorname{ht}(V)$ of a subspace
$V\subseteq\Q^N$ is the largest bit length of an entry of its reduced row echelon basis.
Condition~(3) of Theorem~\ref{thm:main} is the conjunction of
\begin{itemize}
\item[(a)] $\cG^\circ$ is triangularizable over $\Q$, and
\item[(b)] any two unipotent elements of $\cG$ commute.
\end{itemize}

\begin{theorem}\label{thm:polytime}
\begin{enumerate}
\item \emph{(Quasi-unipotent sector.)} Given a rational weighted automaton, one can decide in
deterministic polynomial time whether its series $f$ is invertible and quasi-unipotent, and if
so whether $f$ is CCRA-computable, equivalently whether $\Gamma_f$ is virtually abelian. The
algorithm uses only exact linear algebra over $\Q$ and Hermite normal forms of integer matrices.
\item \emph{(Invertible sector.)} CCRA-computability of an invertible series is decided by a
deterministic algorithm whose running time is polynomial in the input length and in the number
$B_f$ of Definition~\ref{def:heights} (the largest height of the Lie algebras of the unipotent
radicals of those truncations of $\cG_f$ that the algorithm visits), and is $2^{O(m\log(dk))}$ times a polynomial in the input
length, where $m\le d$ is the nilpotency index of the radical of $\spn\Gamma_f$. In
particular it runs in polynomial time when $m$ is bounded, and in the quasi-unipotent sector. By
Theorem~\ref{thm:detpoly}, $B_f$ is always polynomially bounded, so this algorithm runs in
deterministic polynomial time on all invertible series.
\item \emph{(Randomized.)} CCRA-computability of invertible series is decidable in Monte Carlo
polynomial time (bounded two-sided error).
\end{enumerate}
The algorithms in (2) and (3) use the deterministic polynomial-time algorithm of Friedl and
R\'onyai~\cite{FR85} (see also~\cite[Theorem~2.6]{Bea99}) for the simple components of a
semisimple algebra over $\Q$, and rational root finding for polynomials over $\Q$.
\end{theorem}

Theorem~\ref{thm:detpoly} (Section~\ref{ssec:detpoly}) shows that $B_f$ is always bounded by a
polynomial in the input length, extending Lemma~\ref{lem:apriori} from the quasi-unipotent sector.
With Theorem~\ref{thm:polytime}(2) this gives a deterministic polynomial-time algorithm for the
whole invertible sector, so part~(3) is superseded; we keep it because its proof is short.

The proof has three ingredients: a block-triangular form whose diagonal part is completely
reducible (Section~\ref{ssec:block}); a test of~(a) by finiteness testing
(Section~\ref{ssec:atest}); and a layer-by-layer computation of the Lie algebra of the unipotent
radical in which the possibly exponential finite quotient $\cG/\cG^\circ$ is handled implicitly,
through commutants and a cohomological minimality argument, and never enumerated
(Section~\ref{ssec:layers}). Section~\ref{ssec:detpoly} then bounds the heights of the computed Lie
algebras by a polynomial. In outline, the deterministic polynomial-time claim for the invertible
sector is proved as follows: Lemma~\ref{lem:atest} tests~(a); Proposition~\ref{prop:layer}
computes the Lie algebra of $U$ one distance layer at a time; the size analysis
(Lemma~\ref{lem:sizes} and the proof of Theorem~\ref{thm:polytime}(2))
bounds the running time polynomially in the input length and in $B_f$; and
Theorem~\ref{thm:detpoly}, via Lemmas~\ref{lem:iter}--\ref{lem:Upp} and
Proposition~\ref{prop:ugen}, bounds $B_f$ by a polynomial.

\subsection{Block form and the reductive quotient}\label{ssec:block}
Let $A=\spn\Gamma$, computed by closing $\{I\}$ under right multiplication by the $M_a$
(Section~\ref{sec:algorithms}); it is the span of the words of length less than $d^2$, and
$M_a^{-1}\in A$. In characteristic zero its radical is
$J=\{x\in A:\tr(xy)=0\text{ for all }y\in A\}$: this is an ideal consisting of elements $x$ with
$\tr(x^i)=0$ for all $i\ge1$, hence a nil ideal, and every element of the radical satisfies the
trace condition. Put $F_i=J^i\Q^d$ and let $m$ be the least index with $F_m=0$. The $F_i$ are
$A$-submodules, and $F_i/F_{i+1}$ is annihilated by $J$, hence a module over the semisimple
algebra $A/J$ and therefore semisimple.

Choose subspaces $E_p$ with $F_{m-1-p}=E_p\oplus F_{m-p}$ ($0\le p<m$) and a basis of $\Q^d$
adapted to $\Q^d=E_0\oplus\cdots\oplus E_{m-1}$ (computable in polynomial time, with polynomial
bit size). In this basis $\Gamma$ lies in the group $P$ of invertible block upper triangular
matrices. Let $L_P$ be the block diagonal matrices, $U_P$ the block unipotent ones, and
$b:P\to L_P$ the homomorphism that keeps the diagonal blocks. For $t\ge0$ let $\mathfrak v_t$ be the
space of matrices supported on the blocks $(q,p)$ with $p-q=t$, the \emph{distance-$t$ part};
$\fu_P=\bigoplus_{t\ge1}\mathfrak v_t$ is the Lie algebra of $U_P$, and
$X\mapsto X_t$ denotes the projection onto $\mathfrak v_t$. Matrices supported at distance
$>j$ form an ideal $\mathfrak I_{j+1}$ of the algebra of block upper triangular matrices; let
$\pi_j$ be reduction modulo $\mathfrak I_{j+1}$ (delete all blocks at distance $>j$), so that
$P_j=\pi_j(P)$ is again a linear algebraic group and $\pi_j:P\to P_j$ a morphism defined over
$\Q$. We write $\fu_{P,j}=\pi_j(\fu_P)=\bigoplus_{1\le t\le j}\mathfrak v_t$. Note that
$\mathfrak v_{j+1}$ is central in the Lie algebra $\fu_{P,j+1}$, and that $\operatorname{Ad}(u)$
acts trivially on $\mathfrak v_{j+1}$ for $u\in\pi_{j+1}(U_P)$.

Put $\cG_j=\pi_j(\cG)$, the closure of $\Gamma_j=\pi_j(\Gamma)$ (A2), let $H$ be the closure of
$b(\Gamma)$, and write $b_a=b(M_a)$. Since $b$ factors through every $\pi_j$, $b(\cG_j)=H$.

\begin{lemma}\label{lem:radical}
$U:=\cG\cap U_P$ is the unipotent radical $R_u(\cG)$, and for every $j$,
$U_j:=\pi_j(U)=\cG_j\cap\pi_j(U_P)=R_u(\cG_j)$. In particular $U=U_{m-1}$.
\end{lemma}
\begin{proof}
A subspace of $E_p$ is $H$-stable iff it is $b(\Gamma)$-stable, and $E_p\cong F_{m-1-p}/F_{m-p}$ as
$b(\Gamma)$-modules, so $H$ acts completely reducibly on each $E_p$. Let $N=b(R_u(\cG_j))$, a connected
unipotent normal subgroup of $H$ (A2). For an irreducible $H$-submodule $W$ of some $E_p$, the fixed
space of $N$ on $W\otimes\Qb$ is nonzero (Kolchin) and, as $N$ is defined over $\Q$, Galois-stable,
hence spanned by its rational points; so $W^N\ne0$, and it is $H$-stable, so $W^N=W$; hence $N$ acts trivially on
$\bigoplus_pE_p$ and $N=1$, i.e.\ $R_u(\cG_j)\subseteq\ker(b|_{\cG_j})=\cG_j\cap\pi_j(U_P)$.
Conversely $\cG_j\cap\pi_j(U_P)$ is a closed normal unipotent subgroup, connected in
characteristic zero, hence contained in $R_u(\cG_j)$. Finally, if $x=\pi_j(g)\in\cG_j\cap\pi_j(U_P)$
with $g\in\cG$, then $b(g)=b(x)=1$, so $g\in U$; thus $\pi_j(U)=\cG_j\cap\pi_j(U_P)$. For $j=m-1$,
$\pi_j$ is the identity.
\end{proof}

\begin{lemma}\label{lem:acrit}
Condition~(a) holds iff $H^\circ$ is a torus that is split over $\Q$. If (a) holds, then (b)
holds iff $U$ is commutative; thus condition~(3) is equivalent to: $H^\circ$ is a $\Q$-split torus
and $U$ is commutative.
\end{lemma}
\begin{proof}
If $H^\circ$ is a $\Q$-split torus, it is diagonalizable over $\Q$ on each $E_p$ (A5); refining the
block basis by rational weight bases shows that $\cG^\circ\subseteq P$, whose diagonal blocks lie in
$b(\cG^\circ)=H^\circ$ (A2), is upper triangular. Conversely, if $\cG^\circ$ fixes a complete rational
flag, then $H^\circ=b(\cG^\circ)$ fixes the induced flags of the $E_p\cong F_{m-1-p}/F_{m-p}$, so
$H^\circ$ is connected, solvable and, by the proof of Lemma~\ref{lem:radical}, has trivial unipotent
radical; it is therefore a torus (A5), and it is split over $\Q$ because it is triangularizable over
$\Q$ (A5). If (a) holds, every unipotent element of $\cG$ lies in $\cG^\circ$ (A3), and the unipotent
elements of the connected solvable group $\cG^\circ$ form the connected normal subgroup
$(\cG^\circ)_u$ (A5), which equals $R_u(\cG^\circ)=R_u(\cG)=U$ by Lemma~\ref{lem:radical}.
\end{proof}

\subsection{Testing condition (a)}\label{ssec:atest}
\begin{lemma}[Finiteness]\label{lem:finite}
Let $\Delta=\langle g_1,\ldots,g_k\rangle\le\GL_n(\Q)$ be given by its generators. Then:
\begin{enumerate}
\item finiteness of $\Delta$ is decidable in deterministic polynomial time;
\item if $\Delta$ is finite, then $|\Delta|\le(2n+1)^{n^2}$, and every element of $\Delta$ has entries
of bit length polynomial in $n$ and the input length; more precisely, if $w_1,\ldots,w_r$ is a basis
of $\spn\Delta$ consisting of positive words of length $<r$ in the generators, and
$w_1^\#,\ldots,w_r^\#$ is the dual basis for the trace form, then
$h=\sum_i\tr(hw_i)\,w_i^\#$ with $\tr(hw_i)\in\{-n,\ldots,n\}$ for every $h\in\Delta$.
\end{enumerate}
\end{lemma}
Deterministic polynomial-time finiteness testing over $\Q$ is due to Babai, Beals and
Rockmore~\cite{BBR93}; we include a short proof of the version we need.
\begin{proof}
Compute words $w_1=I,\ldots,w_r$ forming a basis of $\mathcal A=\spn\Delta$ by span closure and
the Gram matrix $T=(\tr(w_iw_j))_{ij}$. If $\Delta$ is finite, then: every element has finite order,
so its trace is a sum of $n$ roots of unity lying in $\Q$, hence an integer of absolute value
$\le n$; $\mathcal A$ is semisimple (Maschke), so the trace form is nondegenerate (its radical is a nil
ideal, as in Section~\ref{ssec:block}) and $T$ is invertible. Therefore (2) follows by expanding $h$
in the dual basis, since $hw_i\in\Delta$; the entries of $h$ are integer combinations, with
coefficients of absolute value at most $n$, of the entries of the $w_i^\#=\sum_j(T^{-1})_{ij}w_j$,
which have polynomial bit length, and $h\mapsto(\tr(hw_i))_i$ is injective with values in
$\{-n,\ldots,n\}^r$, $r\le n^2$.

The test rejects (declares $\Delta$ infinite) if $T$ is singular or has an entry outside
$\{-n,\ldots,n\}$. Otherwise let $\Lambda^\#=\{y\in\mathcal A:\tr(yw_i)\in\Z\ \forall i\}$, a lattice
containing $\Lambda_0=\bigoplus_i\Z w_i$ with index $|\det T|$. Compute
$\Lambda_{s+1}=\Lambda_s+\sum_a\Lambda_sg_a$ (Hermite normal forms of integer coordinate matrices
with respect to the dual basis~\cite{KB79}); reject as soon as some $\Lambda_s\not\subseteq\Lambda^\#$.
If $\Delta$ is finite, every $\Lambda_s$ is spanned by positive words, i.e.\ elements of $\Delta$, so
$\Lambda_s\subseteq\Lambda^\#$. Each strict inclusion at least doubles the index over $\Lambda_0$, so
after at most $\log_2|\det T|+1$ steps either we have rejected or $\Lambda:=\Lambda_s$ satisfies
$\Lambda g_a\subseteq\Lambda$ for all $a$. Then $\Lambda$ contains all positive words, and it is spanned
by positive words, so it is the ring $\mathcal O$ spanned by them, a finitely generated $\Z$-module;
$\mathcal O\Z^n$ is a lattice in $\Q^n$ with $g_a\mathcal O\Z^n\subseteq\mathcal O\Z^n$.

Finally let $\rho(g)X=g^{\mathsf T}Xg$ on symmetric matrices $\operatorname{Sym}_n(\Q)$, let
$S^\Delta=\bigcap_a\ker(\rho(g_a)-1)$ and $R=\sum_a\operatorname{im}(\rho(g_a)-1)$. Reject unless
$\operatorname{Sym}_n=S^\Delta\oplus R$; then write $I=I_0+I_1$ with $I_0\in S^\Delta$, $I_1\in R$, and
reject unless $I_0$ is positive definite (leading principal minors). Note that
$\rho(gh)=\rho(h)\rho(g)$, so $\rho(gh)-1=(\rho(h)-1)\rho(g)+(\rho(g)-1)$ and
$\rho(g^{-1})-1=-(\rho(g)-1)\rho(g^{-1})$; hence $R=\sum_{h\in\Delta}\operatorname{im}(\rho(h)-1)$.
If $\Delta$ is finite, the averaging operator $\alpha=\frac1{|\Delta|}\sum_{h\in\Delta}\rho(h)$
satisfies $\alpha(\rho(g)-1)=0$, fixes $S^\Delta$ pointwise and has image in $S^\Delta$, and
$X-\alpha X=\frac1{|\Delta|}\sum_h(1-\rho(h))X\in R$; so $\operatorname{Sym}_n=S^\Delta\oplus R$,
$\alpha$ is the projection onto $S^\Delta$ along $R$, and $I_0=\alpha(I)=\frac1{|\Delta|}\sum_hh^{\mathsf T}h$
is positive definite. Hence finite groups are accepted.
Conversely, if all tests pass, $\Delta$ preserves the positive definite form $I_0$, so
$|\det g_a|=1$ and therefore $g_a\mathcal O\Z^n=\mathcal O\Z^n$; thus $\Delta$ lies in the
intersection of the discrete group $\GL(\mathcal O\Z^n)$ with the compact group $O(I_0)(\mathbb R)$,
which is finite. All steps are polynomial: the words have length $<n^2$, the Hermite normal forms
involve integers bounded by $|\det T|\le(n\sqrt r)^r$, and the last step is linear algebra in
dimension $n(n+1)/2$.
\end{proof}

\begin{lemma}\label{lem:atest}
Let $\theta(g)X=gXg^{-1}$ be the conjugation action of $b(\Gamma)$ on $M_d(\Q)$ and
$A_\theta=\spn\theta(b(\Gamma))$. Condition (a) holds iff $A_\theta$ is semisimple and the following
algebra $C$ is commutative and the minimal polynomial of every element of some (equivalently, every)
basis of $C$ splits into linear factors over $\Q$: $C$ is the centralizer in $M_d(\Q)$ of $\mathcal B$,
the sum of those isotypic components $eM_d(\Q)$ ($e$ a primitive central idempotent of $A_\theta$) on
which $\theta(b(\Gamma))$ acts as a finite group. Consequently (a) is decidable in deterministic
polynomial time.
\end{lemma}
\begin{proof}
Let $\Q[H^\circ]$ denote the span of $b(\Gamma)\cap H^\circ$, which is Zariski dense in $H^\circ$
(Lemma~\ref{lem:monoid}); its $\Qb$-span is the span of $H^\circ$. If $C$ is commutative and a basis
$c_1,\ldots,c_s$ has split minimal polynomials, then every element of $C$ has only rational eigenvalues:
modulo the radical, $C$ is a product of number fields, in each of which the images of the $c_i$ are
roots of their minimal polynomials, hence rational, and they span; so each field is $\Q$.

Suppose the conditions hold. Each $eM_d(\Q)$ on which the action is finite is fixed pointwise by
$H^\circ$, because the closure of a finite group is finite and $H^\circ$ is connected; hence
$\mathcal B\subseteq M_d(\Q)^{H^\circ}$, which is the centralizer of $\Q[H^\circ]$, and so
$\Q[H^\circ]\subseteq C$. Thus $H^\circ$ is commutative. Its unipotent elements form a closed connected
subgroup~\cite[4.7]{Bor91}, normal in $H$, which is trivial by the argument in the proof of
Lemma~\ref{lem:radical}; so all elements of $H^\circ$ are semisimple and $H^\circ$ is a torus (A4). The elements
of $b(\Gamma)\cap H^\circ\subseteq C$ are therefore commuting semisimple matrices with rational
eigenvalues, hence simultaneously diagonalizable over $\Q$, and so is their closure $H^\circ$. Thus
$H^\circ$ is a $\Q$-split torus, and (a) holds by Lemma~\ref{lem:acrit}.

Conversely suppose (a), so $H^\circ$ is a $\Q$-split torus and $H/H^\circ$ is finite. Every rational
$H$-module is then completely reducible (decompose under $H^\circ$ into weight spaces, take an
$H^\circ$-stable complement and average the projection over $H/H^\circ$), so $M_d(\Q)$ is a semisimple
$A_\theta$-module, $A_\theta$ is semisimple, and $M_d(\Q)=\bigoplus_eeM_d(\Q)$ is the isotypic
decomposition for $b(\Gamma)$, equivalently for $H$. The subspace $M_d(\Q)^{H^\circ}$ is $H$-stable;
on an irreducible $H$-submodule $V$ the fixed space $V^{H^\circ}$ is $0$ or $V$, so each isotypic
component lies in $M_d(\Q)^{H^\circ}$ or meets it trivially, and $M_d(\Q)^{H^\circ}$ is the sum of the
components it contains. On these $H$ acts through the finite group $H/H^\circ$, while a component
on which $\theta(b(\Gamma))$ is finite is fixed by $H^\circ$ as above. Hence
$\mathcal B=M_d(\Q)^{H^\circ}$, the centralizer of the commutative split semisimple algebra
$\Q[H^\circ]=\bigoplus_\chi\Q\varepsilon_\chi$ (weight idempotents). This centralizer is
$\bigoplus_\chi\operatorname{End}(\varepsilon_\chi\Q^d)$, whose centralizer is
$\bigoplus_\chi\Q\varepsilon_\chi$. So $C=\Q[H^\circ]$, which is commutative and diagonal over $\Q$.

Algorithm: compute $A_\theta\subseteq\operatorname{End}(M_d(\Q))$ by span closure, test
semisimplicity (radical by the trace criterion), compute its primitive central idempotents by the
algorithm of Friedl and R\'onyai~\cite{FR85}, test finiteness on each component by
Lemma~\ref{lem:finite}, compute $C$ by linear algebra, test commutativity, and test that the minimal
polynomial of each basis element of $C$ splits into linear factors over $\Q$ (find its rational roots,
e.g.\ by factorization over $\Q$, and compare degrees). All matrices involved have size at most
$d^2\times d^2$ and polynomial bit length. The semisimplicity test in fact always succeeds. Indeed, $H$ acts
completely reducibly on $\bigoplus_pE_p$ over $\Q$, hence over $\Qb$ (characteristic zero); the argument
in the proof of Lemma~\ref{lem:radical}, applied to the unipotent radical $R_u(H)$ itself, shows that
$R_u(H)$ acts trivially on every irreducible $H$-submodule, so $R_u(H)=1$ and $H$ is reductive; and for
a reductive group in characteristic zero (averaging over the finite group $H/H^\circ$ when $H$ is
disconnected) the module $M_d(\Qb)\cong\Qb^d\otimes(\Qb^d)^*$ is completely reducible, so $A_\theta$ is
semisimple. Thus (a) holds iff $C$ is commutative and split; we keep the semisimplicity condition in
the statement because the primitive central idempotents of $A_\theta$ are needed to define
$\mathcal B$.
\end{proof}

In the quasi-unipotent sector no idempotents are needed:

\begin{lemma}\label{lem:qutest}
An invertible series is quasi-unipotent iff $b(\Gamma)$ is finite. In this case (a) holds and
$U=\cG^\circ$, so $f$ is CCRA-computable iff $U$ is commutative.
\end{lemma}
\begin{proof}
Word matrices and their block diagonal parts have the same eigenvalues. If $b(\Gamma)$ is finite, all
eigenvalues of word matrices are roots of unity. Conversely, if $f$ is quasi-unipotent then by the
proof of Lemma~\ref{lem:qu} all elements of $\cG$, hence of $H=b(\cG)$, satisfy $(g^N-I)^d=0$ for
some $N$; so $H^\circ$ is unipotent (A4), and it is trivial because $H$ acts completely reducibly
(proof of Lemma~\ref{lem:radical}); thus $H$ and $b(\Gamma)$ are finite. In that case $H^\circ=1$
is a split torus, (a) holds, $\cG^\circ$ is unipotent (Lemma~\ref{lem:qu}) and equals
$R_u(\cG)=U$; conclude by Lemma~\ref{lem:acrit}.
\end{proof}

\subsection{The unipotent radical, layer by layer}\label{ssec:layers}
From now on assume (a), so that $H^\circ$ is a $\Q$-split torus and every rational $H$-module, over $\Q$
or over $\Qb$, is completely reducible (proof of Lemma~\ref{lem:atest}; over $\Q$ the coset
representatives of $H/H^\circ$ can be taken in $b(\Gamma)$, which meets every coset).

\begin{lemma}[Abelian extensions]\label{lem:abext}
Let $E=L\ltimes\mathcal M$ be a linear algebraic group over $\Q$, where $\mathcal M$ is a vector
group, identified with its Lie algebra and acted on by $L$ through a rational representation
$\operatorname{Ad}$, and every rational $L$-module, over $\Q$ or over $\Qb$, is completely reducible. Let
$g_a=\ell_a\exp(c_a)\in E(\Q)$ $(1\le a\le k)$ with $\ell_a\in L(\Q)$ generating a Zariski dense
subgroup of $L$ and $c_a\in\mathcal M(\Q)$, and let $\mathcal D$ be the closure of
$\langle g_1,\ldots,g_k\rangle$. Let $\mathcal E$ be the commutant of
$\{\operatorname{Ad}(\ell_a)|_{\mathcal M(\Q)}\}_a$ and
\[
\Phi=\bigl\{\varphi\in\mathcal E:\ \exists y\in\mathcal M(\Q)\ \forall a\quad
\varphi(c_a)=(1-\operatorname{Ad}(\ell_a)^{-1})\,y\bigr\}.
\]
Then $\mathcal D\cap\mathcal M=\exp W$ for the subspace $W=\bigcap_{\varphi\in\Phi}\ker\varphi$ of
$\mathcal M(\Q)$, and there is $y\in\mathcal M(\Q)$ with
$c_a-(1-\operatorname{Ad}(\ell_a)^{-1})y\in W$ for all $a$.
\end{lemma}
\begin{proof}
For $g\in\Gamma'=\langle g_a\rangle$ write $g=\ell(g)\exp(c(g))$; then $\ell$ is a homomorphism and
$c(gh)=\operatorname{Ad}(\ell(h))^{-1}c(g)+c(h)$, so $c$ is a $1$-cocycle for the right action
$X\cdot g=\operatorname{Ad}(\ell(g))^{-1}X$, determined by its values $c_a$ on generators (and
$c(g^{-1})=-\operatorname{Ad}(\ell(g))c(g)$). Call a subspace $W'\subseteq\mathcal M(\Q)$ \emph{working}
if it is $\operatorname{Ad}(L)$-stable and some $y$ satisfies
$c_a\equiv(1-\operatorname{Ad}(\ell_a)^{-1})y\pmod{W'}$ for all $a$; since both sides are cocycles,
the congruence then holds for all $g\in\Gamma'$.

$\mathcal D\cap\mathcal M$ is a closed subgroup of a vector group, hence a subspace $\exp W_0$ defined
over $\Q$. It is normalized by $\mathcal D$, and since $\mathcal M$ is commutative, conjugation by
$\ell\exp(X)$ acts on $\mathcal M$ as $\operatorname{Ad}(\ell)$; the projection of $\mathcal D$ to $L$ is
closed and contains $\langle\ell_a\rangle$, hence is $L$, so $W_0$ is $\operatorname{Ad}(L)$-stable.

\emph{$W_0$ is working.} The image of $\mathcal D$ in $E/\exp W_0=L\ltimes(\mathcal M/W_0)$ meets
$\mathcal M/W_0$ trivially and maps onto $L$ bijectively, hence isomorphically (characteristic zero);
so it is the graph $\{\ell\exp s(\ell)\}$ of a morphism $s:L\to\mathcal M/W_0$ satisfying
$s(\ell\ell')=\operatorname{Ad}(\ell')^{-1}s(\ell)+s(\ell')$. On $\mathcal M/W_0\oplus\Qb$ let
$\ell$ act by $(x,\tau)\mapsto(\operatorname{Ad}(\ell)x-\tau\operatorname{Ad}(\ell)s(\ell),\tau)$,
a rational representation (the cocycle identity gives the homomorphism property). The
submodule $\mathcal M/W_0\oplus0$ has an $L$-stable complement, spanned by some $(y_0,1)$; then
$\operatorname{Ad}(\ell)y_0-\operatorname{Ad}(\ell)s(\ell)=y_0$, i.e.\
$s(\ell)=(1-\operatorname{Ad}(\ell)^{-1})y_0$. Hence
$c_a\equiv s(\ell_a)=(1-\operatorname{Ad}(\ell_a)^{-1})y_0$ modulo $W_0$. This is a linear system with
rational coefficients that has a solution over $\Qb$, so it has one over $\Q$.

\emph{Every working $W'$ contains $W_0$.} If $y$ witnesses that $W'$ is working, then, because
$\mathcal M$ is commutative,
$\exp(y)\,g\exp(-y)=\ell(g)\exp\bigl(c(g)-(1-\operatorname{Ad}(\ell(g))^{-1})y\bigr)\in L\ltimes\exp W'$
for all $g\in\Gamma'$. The latter is a closed subgroup, so it contains $\exp(y)\mathcal D\exp(-y)$,
and intersecting with $\mathcal M$ (on which conjugation by $\exp(y)$ is trivial) gives
$W_0\subseteq W'$.

\emph{$W=W_0$.} The commutant $\mathcal E$ is the algebra of $L$-equivariant endomorphisms of the
completely reducible module $\mathcal M(\Q)$, a semisimple algebra, and $\Phi$ is a left ideal of
$\mathcal E$: if $\varphi(c_a)=(1-\operatorname{Ad}(\ell_a)^{-1})y$ then
$(\psi\varphi)(c_a)=(1-\operatorname{Ad}(\ell_a)^{-1})\psi y$ for $\psi\in\mathcal E$. Hence
$\Phi=\mathcal E\pi_0$ for an idempotent $\pi_0\in\Phi$, and $W=\ker\pi_0$. Since
$\mathcal M(\Q)=\ker\pi_0\oplus\operatorname{im}\pi_0$ as $L$-modules, $\pi_0\in\Phi$ says exactly that
$\ker\pi_0$ is working; so $W_0\subseteq W$. Conversely let $\pi$ be the $L$-equivariant projection
onto an $L$-stable complement of $W_0$ along $W_0$; applying $\pi$ to the congruence that makes $W_0$
working shows $\pi\in\Phi$, so $W\subseteq\ker\pi=W_0$.
\end{proof}

Call $v\in U_P(\Q)$ a \emph{Levi conjugator at level $j$} if $\pi_j(vb_av^{-1})\in\cG_j$ for all $a$,
equivalently $\pi_j(vHv^{-1})\subseteq\cG_j$, equivalently
$\pi_j\bigl((vb_av^{-1})^{-1}M_a\bigr)\in U_j$ for all $a$. The identity is a Levi conjugator at level
$0$. Let $\fu_j\subseteq\fu_{P,j}$ be the Lie algebra of $U_j$; when $U_j$ is commutative it is a vector
group and $U_j=\exp\fu_j$ (this is the only case used).

\begin{proposition}[Layer step]\label{prop:layer}
Assume (a). Let $0\le j\le m-2$, assume $U_j$ is commutative, and let $v$ be a Levi conjugator at
level $j$. View elements of $\fu_{P,j}$ as elements of $\fu_{P,j+1}$ with zero distance-$(j+1)$ part.
\begin{enumerate}
\item $U_{j+1}$ is commutative iff $[X,Y]_{j+1}=0$ for all $X,Y\in\fu_j$.
\item Suppose it is. Put $\mathcal M=\fu_j\oplus\mathfrak v_{j+1}$,
$\ell_a=\pi_{j+1}(vb_av^{-1})$ and $c_a=\log\bigl(\ell_a^{-1}\pi_{j+1}(M_a)\bigr)$. Then
$\mathcal M$ is a commutative Lie algebra stable under every $\operatorname{Ad}(\ell_a)$, each $c_a$
lies in $\mathcal M$, and Lemma~\ref{lem:abext}, applied with $L=\pi_{j+1}(vHv^{-1})$, these
$\ell_a$, $c_a$, yields $W=\fu_{j+1}$.
\item There is $y\in\mathfrak v_{j+1}(\Q)$ with $c_a-(1-\operatorname{Ad}(b_a)^{-1})y\in\fu_{j+1}$ for
all $a$, and for every such $y$ the element $\exp(-y)v$ is a Levi conjugator at level $j+1$.
\end{enumerate}
\end{proposition}
\begin{proof}
(1) $\mathcal M$ is the preimage of $\fu_j$ in $\fu_{P,j+1}$, and $\fu_{j+1}\subseteq\mathcal M$ maps onto
$\fu_j$. Since $\mathfrak v_{j+1}$ is central and $[X,Y]\in\mathfrak v_{j+1}$ for $X,Y\in\fu_j$ (the
image of $[X,Y]$ in $\fu_{P,j}$ vanishes), $[\mathcal M,\mathcal M]$ is spanned by the $[X,Y]_{j+1}$. If
these vanish, $\fu_{j+1}\subseteq\mathcal M$ is commutative. Conversely, lifts $\tilde X,\tilde Y\in\fu_{j+1}$
of $X,Y$ differ from $X,Y$ by central elements, so $[X,Y]_{j+1}=[\tilde X,\tilde Y]=0$ if $\fu_{j+1}$ is
commutative.

(2) Let $\mathcal E_j\subseteq P_{j+1}$ be the preimage of $\cG_j$; then
$\tilde U:=\mathcal E_j\cap\ker b$ is the preimage of $U_j$, with Lie algebra $\mathcal M$, a vector
group by (1). As $v$ is a Levi conjugator, $L=\pi_{j+1}(vHv^{-1})\subseteq\mathcal E_j$; $b$ maps $L$
isomorphically onto $H$ and kills $\tilde U$, so $\mathcal E_j=L\ltimes\tilde U$, and rational $L$-modules
are completely reducible. We have $\ell_a\in L$, $b(\ell_a)=b_a$, so the $\ell_a$ generate a dense
subgroup of $L$; $\ell_a^{-1}\pi_{j+1}(M_a)\in\mathcal E_j\cap\ker b=\tilde U$ gives $c_a\in\mathcal M$.
Now $\cG_{j+1}\subseteq\mathcal E_j$ is the closure of $\langle\pi_{j+1}(M_a)\rangle$ and
$\cG_{j+1}\cap\tilde U=\cG_{j+1}\cap\ker b=U_{j+1}$ (Lemma~\ref{lem:radical}). Lemma~\ref{lem:abext}
gives $\exp W=U_{j+1}$, i.e.\ $W=\fu_{j+1}$.

(3) Lemma~\ref{lem:abext} gives $y_0\in\mathcal M(\Q)$ with
$c_a-(1-\operatorname{Ad}(\ell_a)^{-1})y_0\in\fu_{j+1}$. Write $y_0=w+y$ with $w\in\fu_{j+1}$,
$y\in\mathfrak v_{j+1}$ (possible as $\fu_{j+1}+\mathfrak v_{j+1}=\mathcal M$). Conjugation by any
element of $\mathcal E_j$ acts on $\mathcal M$ as conjugation by an element of $\cG_{j+1}$ with the same
image under $b$, so it preserves $\fu_{j+1}$; hence $y$ also works, and on $\mathfrak v_{j+1}$,
$\operatorname{Ad}(\ell_a)=\operatorname{Ad}(b_a)$ because $\operatorname{Ad}(\pi_{j+1}(v))$ is trivial there.
Then, $\mathcal M$ being commutative, $\exp(y)\pi_{j+1}(M_a)\exp(-y)
=\ell_a\exp\bigl(c_a-(1-\operatorname{Ad}(\ell_a)^{-1})y\bigr)\in\ell_aU_{j+1}$, i.e.\
$\pi_{j+1}\bigl((\exp(-y)vb_av^{-1}\exp(y))^{-1}M_a\bigr)\in\exp(-y)U_{j+1}\exp(y)=U_{j+1}$.
\end{proof}

\begin{definition}\label{def:heights}
For an invertible series $f$ satisfying (a), let $j^*$ be the largest $j\le m-1$ such that
$U_0,\ldots,U_{j^*}$ are commutative, and $B_f=\max_{j\le j^*}\operatorname{ht}(\fu_j)$, heights taken in
the adapted basis with the matrix coordinates ordered by increasing distance. (Put $B_f=0$ if (a)
fails.) $B_f$ depends on the choice of the adapted basis; the algorithm uses the basis computed in step~1.
\end{definition}

\paragraph{The algorithm.}
\begin{enumerate}
\item Minimize, test invertibility (Theorem~\ref{thm:decision}); compute $A$, $J$, the flag, the
adapted basis, $b_a$.
\item Test (a) by Lemma~\ref{lem:atest}; if (a) fails, answer NO. (For part~(1) of
Theorem~\ref{thm:polytime} one instead tests finiteness of $b(\Gamma)$ by Lemma~\ref{lem:finite}; if
$b(\Gamma)$ is infinite, $f$ is not quasi-unipotent (Lemma~\ref{lem:qutest}) and we report this; otherwise
(a) holds and we proceed.)
\item Put $\fu_0=0$, $v=I$. For $j=0,\ldots,m-2$: if $[X,Y]_{j+1}\ne0$ for two basis elements of
$\fu_j$, answer NO. Otherwise compute $\ell_a,c_a$, the matrices of $\operatorname{Ad}(\ell_a)$ on
$\mathcal M$, the commutant $\mathcal E$, a basis $\varphi_1,\ldots,\varphi_s$ of $\Phi$ (projection of
the solution space of a homogeneous linear system in $(\varphi,y)$), and
$\fu_{j+1}=\bigcap_i\ker\varphi_i$, stored by its reduced echelon basis with respect to coordinates
ordered by increasing distance. Then compute the new Levi conjugator: let $u_a\in\fu_{j+1}$ be the
combination of the echelon basis vectors whose pivots have distance $\le j$, with coefficients the
entries of $c_a$ at these pivots (so $u_a\equiv c_a$ modulo $\mathfrak v_{j+1}$), let
$r_a=(c_a-u_a)_{j+1}$, solve
\[
(1-\operatorname{Ad}(b_a)^{-1})\,y+w_a=r_a\qquad(1\le a\le k),\quad y\in\mathfrak v_{j+1},\
w_a\in\fu_{j+1}\cap\mathfrak v_{j+1},
\]
with free variables set to zero, and replace $v$ by $\pi_{j+1}(\exp(-y)v)$. (The basis vectors with pivot
at distance $j+1$ are supported at distance $j+1$ and span $\fu_{j+1}\cap\mathfrak v_{j+1}$; the solutions
$y$ of this system are exactly the $y\in\mathfrak v_{j+1}$ of Proposition~\ref{prop:layer}(3).)
\item Answer YES.
\end{enumerate}
By Lemmas~\ref{lem:radical}, \ref{lem:acrit}, \ref{lem:qutest} and Proposition~\ref{prop:layer}, the
answer is YES iff condition (3) holds (note $U_0=1$ and $U=U_{m-1}$), i.e.\ iff $f$ is
CCRA-computable. The number of arithmetic operations is polynomial in $d$ and $k$:
$\dim\mathcal M\le d^2$, $\mathcal E$ is the kernel of a linear map on $d^4$ unknowns, and $m\le d$.

\subsection{Bit complexity}\label{ssec:bits}
\begin{lemma}[A priori heights in the quasi-unipotent sector]\label{lem:apriori}
If $b(\Gamma)$ is finite, then $\operatorname{ht}(\fu_j)$ is bounded by a polynomial in $d$ and $\sigma$ for every
$j$ (in the adapted basis, whose bit size is polynomial).
\end{lemma}
\begin{proof}
Let $K=\Gamma\cap U_P$, the kernel of $b|_\Gamma$, of index $|b(\Gamma)|\le(2d+1)^{d^2}$
(Lemma~\ref{lem:finite}). By Schreier's lemma $K$ is generated by elements $r_hM_ar_{hb_a}^{-1}$, where
$r_h$ is a shortest positive word with $b(r_h)=h\in b(\Gamma)$; these are words of length
$<2|b(\Gamma)|$ in the $M_a^{\pm1}$. A word $g=M_{a_1}^{\eps_1}\cdots M_{a_s}^{\eps_s}$ factors as
$g=b(g)\prod_{i=1}^s(1+x_i)$ with $1+x_i=\beta_i^{-1}\,b(M_{a_i}^{\eps_i})^{-1}M_{a_i}^{\eps_i}\,\beta_i$,
$\beta_i=b(M_{a_{i+1}}^{\eps_{i+1}}\cdots M_{a_s}^{\eps_s})\in b(\Gamma)$. By
Lemma~\ref{lem:finite}(2) all elements of $b(\Gamma)$ have polynomial bit length and a common
denominator of polynomial bit length, so the $x_i$ are strictly block upper triangular with entries
in $D^{-1}\Z$ of absolute value $\le2^{p}$, for some $D$ and $p$ polynomial in $d,\sigma$. Since
products of $m$ of the $x_i$ vanish, $\prod_i(1+x_i)$ is a sum of at most $ms^{m-1}$ products of fewer
than $m$ factors $x_i$; hence the entries of $g$ have bit length $O(m\log s+mp+m\log D)$, polynomial
for $s<2|b(\Gamma)|$. The same holds for the logarithms of the Schreier generators. In the quasi-unipotent
sector $K=\Gamma\cap\cG^\circ$ (Lemma~\ref{lem:qutest}) is dense in $U=\cG^\circ$ (Lemma~\ref{lem:monoid}).
The Lie algebra $\mathfrak l$ generated by the logarithms of the Schreier generators equals $\fu$:
$\exp\mathfrak l$ is a closed subgroup (Baker--Campbell--Hausdorff; $\exp$ and $\log$ are inverse
polynomial bijections between $\fu_P$ and $U_P$) containing $K$, so $\fu\subseteq\mathfrak l$; and
$\fu$ contains the logarithm of every unipotent element of $U$ (A3) and is a Lie algebra, so
$\mathfrak l\subseteq\fu$. As $\mathfrak l$ is spanned by left-normed brackets of fewer than $m$ such
logarithms, it, and every $\fu_j=\pi_j(\fu)$, is spanned by vectors of polynomial bit length, and
the reduced echelon basis of such a span has polynomial height.
\end{proof}

\begin{lemma}[Size of the Levi conjugators]\label{lem:sizes}
There is a polynomial $p$ such that the following holds in step~3. If the entries of the $M_a^{\pm1}$
and $b_a^{\pm1}$ (in the adapted basis) and of the echelon bases of $\fu_1,\ldots,\fu_j$ have bit length
at most $\beta\ge1$, then every number computed at level $j+1$, including the echelon basis of
$\fu_{j+1}$, has bit length at most $p(d,k)\,\beta$.
\end{lemma}
\begin{proof}
Write $y_t\in\mathfrak v_t$ for the correction computed at level $t$. Since $y_t^2$ and
$y_t(v-1)$ are supported at distances $>t$, the update $v\mapsto\pi_t(\exp(-y_t)v)$ is
$v\mapsto v-y_t$; hence $v=I-\sum_{t\le j}y_t$ at level $j$. The coefficient matrix $C_t$ of the
system for $y_t$ involves only the $b_a^{\pm1}$ and the echelon basis of $\fu_t\cap\mathfrak v_t$; the
solution with free variables zero is $y_t=\operatorname{adj}(C_t'')\,\gamma_tr'/\det(C_t'')$, where
$C_t''=\gamma_tC_t'$ is an integer matrix, $C_t'$ is a nonsingular square submatrix of $C_t$ determined
by $C_t$ alone, $\gamma_t$ is a common denominator of $C_t'$, and $r'$ collects entries of the $r_a$.
Thus the earlier corrections enter only through the right-hand sides.

Let $\delta$ be the product of $m!$, of $2^{\beta}$, of the denominators of all entries of the
$b_a^{\pm1}$, $M_a^{\pm1}$ and of the echelon bases of $\fu_1,\ldots,\fu_j$, and of the integers
$\gamma_t|\det C_t''|$ for $t\le j$; its bit length is at most $p_1(d,k)\beta$ for a polynomial $p_1$.
For $c,\kappa\ge1$ call a matrix $X$ supported at distances $\ge1$ \emph{$(c,\kappa)$-tame} if, for
every $t\ge1$, $\delta^{ct}X_t$ is an integer matrix and every entry of $X_t$ has absolute value at
most $\kappa^t$. A $(c,\kappa)$-tame matrix is $(c',\kappa')$-tame for $c'\ge c$, $\kappa'\ge\kappa$.
Moreover:
\begin{itemize}
\item if $X$ is $(c,\kappa)$-tame and $Y$ is $(c,\kappa')$-tame, then $XY$ and $X+Y$ are
$(c,dm\max(\kappa,\kappa'))$-tame, since an entry of $(XY)_t$ is a sum of at most $dm$ products of
entries of $X_{t_1}$ and $Y_{t_2}$ with $t_1+t_2=t$;
\item $b_a^{\pm1}Xb_a^{\mp1}$ is $(c+2,d^2\delta^2\kappa)$-tame, since $\delta b_a^{\pm1}$ is an integer
matrix, the entries of $b_a^{\pm1}$ have absolute value at most $2^\beta\le\delta$, and
$\delta^2\le\delta^{2t}$ for $t\ge1$;
\item $Xb_a^{\pm1}$ and $b_a^{\pm1}X$ are $(c+1,d\delta\kappa)$-tame, for the same reasons;
\item $\exp(X)-1$, $\log(1+X)$ and $(1+X)^{-1}-1$ are $(c+1,(dm)^{m+1}\kappa)$-tame, because they
are combinations of $X,\ldots,X^{m-1}$ with coefficients of absolute value at most $1$ of the form
$\pm1/i$ or $1/i!$ with $i<m$, and $i!\mid\delta$;
\item if $X$ is $(c,\kappa)$-tame and $u=\sum_iX_{p_i}e_i$, where the $e_i$ are echelon basis vectors
of some $\fu_t$ with pivot $p_i$ at distance $s_i$ and $X_{p_i}$ is the entry of $X$ at $p_i$, then $u$
is $(c+1,d^2\delta\kappa)$-tame: the entries of $e_i$ vanish at distances $<s_i$ (echelon form with
coordinates ordered by distance), have denominators dividing $\delta$ and absolute value at most
$2^\beta\le\delta$, and $\delta^{cs_i+1}\mid\delta^{(c+1)s}$ for $s\ge s_i\ge1$.
\end{itemize}
Also $\nu_a-1$, where $\nu_a=b_a^{-1}M_a$, is $(2,d\delta^2)$-tame. Suppose $y_1,\ldots,y_{t-1}$ are
$(c,\kappa)$-tame. Then $v-1$ is $(c,dm\kappa)$-tame, and
$\ell_a^{-1}\pi_t(M_a)=\pi_t\bigl(v\cdot b_a^{-1}v^{-1}b_a\cdot\nu_a\bigr)$, $c_a$, $u_a$ and $r_a$ are
obtained from $v-1$ and $\nu_a-1$ by a bounded number of the operations above; hence $r_a$ is
$(c+c_1,F\kappa)$-tame for an absolute constant $c_1$ and an $F$ with $\log F\le p_2(d,k)\beta$,
independent of $t$ and of $c$. Since $|\det C_t''|\ge1$ and $\gamma_t\det C_t''\mid\delta$, $y_t$ is then
$(c+c_1+1,F'\kappa)$-tame with $F'=F\cdot\max_{t'\le m}\bigl(\dim(\mathfrak v_{t'})\,(1+k\dim\mathfrak v_{t'})\cdot\gamma_{t'}
\max|\operatorname{adj}C_{t'}''|\bigr)$ (maximum over the levels at which these systems are defined),
again with $\log F'\le p_3(d,k)\beta$ independent of $t$ and $c$. By
induction on $t$, using monotonicity, every $y_t$ ($t\le j$) is
$(c_0+(c_1+1)t,(F')^t\kappa_0)$-tame for absolute $c_0$ and $\kappa_0=d\delta^2$, and the same holds,
with $t$ replaced by $m$, for $v-1$, $\ell_a^{\pm1}-b_a^{\pm1}$ and $c_a$ at level $j+1$. An entry at
distance $s\le m$ of a $(c,\kappa)$-tame matrix has denominator at most $\delta^{cs}$ and numerator at
most $\kappa^s\delta^{cs}$, hence bit length at most $2cs\log\delta+s\log\kappa+2$; with $c=O(m)$,
$\log\kappa=O(m\,p_3(d,k)\beta)$ and $m\le d$ this is at most $p_4(d,k)\beta$. The matrices of the
$\operatorname{Ad}(\ell_a)$ on $\mathcal M$ (whose basis consists of the echelon basis of $\fu_j$ and elementary
matrices), the commutant $\mathcal E$, the system defining $\Phi$ and the kernel $\bigcap_i\ker\varphi_i$
are then linear-algebra problems of dimension at most $\operatorname{poly}(d,k)$ with coefficients of
bit length $O(p_4(d,k)\beta)$, so their reduced echelon solutions, in particular the echelon basis of
$\fu_{j+1}$, have bit length at most $p_5(d,k)\beta$ (Cramer's rule and Hadamard's bound).

It remains to bound the correction $y_{j+1}$, whose system involves $\fu_{j+1}$. Repeat the argument
with $\beta$ replaced by $\beta'=p_5(d,k)\beta$, which bounds the bit lengths of the echelon bases of
$\fu_1,\ldots,\fu_{j+1}$; this adds the denominators of the echelon basis of $\fu_{j+1}$ and
$\gamma_{j+1}|\det C_{j+1}''|$ to $\delta$ and shows that $y_{j+1}$, and the new $v$, have bit length at most
$p_4(d,k)\beta'$. Taking $p=p_4p_5$ proves the lemma.
\end{proof}

\begin{proof}[Proof of Theorem~\ref{thm:polytime}]
(1) Minimization and invertibility are polynomial (Theorem~\ref{thm:decision}); quasi-unipotence is
decided by Lemma~\ref{lem:qutest} and Lemma~\ref{lem:finite}(1), which uses linear algebra and Hermite
normal forms only. In this sector condition (a) holds and CCRA-computability is equivalent to
commutativity of $U$ (Lemma~\ref{lem:qutest}) and to virtual abelianness of $\Gamma_f$
(Theorem~\ref{thm:sectors}(2)). Step 3 decides commutativity of $U$. By Lemma~\ref{lem:apriori}
every $\fu_j$ has polynomial height $\beta$; since step 3 stores $\fu_{j+1}$ in reduced echelon form,
Lemma~\ref{lem:sizes} applies at every level with the same polynomial $\beta$, so every level runs in
polynomial time, and there are fewer than $d$ levels.

(2) Step 2 is polynomial by Lemma~\ref{lem:atest} and Lemma~\ref{lem:finite}. In the adapted basis the
entries of the $M_a^{\pm1}$ and $b_a^{\pm1}$ have bit length at most $\sigma'$, a polynomial in $d$ and
$\sigma$; below, $\sigma$ stands for $\sigma'$ (the difference is absorbed in the polynomials). Step 3 performs the bracket test only on
$\fu_j$ with $j\le j^*$ and computes $\fu_{j+1}$ only if $j+1\le j^*$, so all heights involved are at most
$B_f$, and so by Lemma~\ref{lem:sizes} its running time is polynomial
in the input length and $B_f$. For the second bound, Lemma~\ref{lem:sizes} and the bit-size bounds of
Gaussian elimination show that the reduced echelon basis of $\fu_{j+1}$ computed at level $j+1$ has
height at most $p(d,k)\,(\sigma+\max_{t\le j}\operatorname{ht}\fu_t)$ for a fixed polynomial $p$. Iterating over at most
$m-1$ levels gives $B_f\le(p(d,k)+1)^{m}(\sigma+1)$, and hence the time bound
$2^{O(m\log(dk))}\operatorname{poly}(\sigma)$.

(3) Run steps 1--2 exactly. By the proof of (2), every rational number occurring in step 3 is a
quotient of integers of bit length at most $\Sigma=2^{O(m\log(dk))}\operatorname{poly}(\sigma)$, and the step
consists of polynomially many arithmetic operations and reduced echelon computations with a fixed
pivoting rule, together with zero tests. All polynomials in Lemma~\ref{lem:sizes} and in the proof of~(2) are explicit (they come from Cramer's rule and Hadamard's bound), so $\Sigma$, and the range $X$ below, can be computed from $d$, $k$, $m$ and $\sigma$. Let $\mathcal Z$ be the set of nonzero integers formed by the
denominators of all rational numbers occurring, the pivot minors of all echelon computations, and the
numerators of all brackets tested in step 3; their product has bit length at most
$\operatorname{poly}(d,k)\,\Sigma$, so at most that many primes divide it. Perform step 3 modulo a prime $\mathfrak q>m$
chosen uniformly among the primes in $[X,2X]$ with $X$ a suitable polynomial in $\Sigma$, so that $\mathfrak q$ has
$O(m\log(dk)+\log\sigma)$ bits. Here $\mathcal Z$ includes, for every echelon computation (also the
projection computing $\Phi$ and the particular solutions), a nonzero minor certifying its rank profile
(nonzero rows on the pivot columns); a prime not dividing it nor any denominator yields the same rank
profile, and the reduced echelon form modulo $\mathfrak q$ is the reduction of the rational one. If
$\mathfrak q$ divides no element of $\mathcal Z$, every echelon form, solution and zero test computed
modulo $\mathfrak q$ is the reduction of the rational one, and the answer is correct; otherwise (for
instance if an inconsistent system is met) the output is arbitrary, which is accounted for in the failure
probability. The random prime is obtained by sampling integers in $[X,2X]$ and testing primality
deterministically. By the prime
number theorem the failure probability is at most $1/3$ for a suitable choice of $X$, and it can be
reduced by independent repetition and majority vote.
\end{proof}

\begin{remark}\label{rem:polytime}
(i) The finite quotient $\cG/\cG^\circ$ may have order exponential in $d$ (for instance the hyperoctahedral
group of order $2^dd!$), and generating sets of $\Gamma\cap\cG^\circ$ obtained from a transversal can be
exponentially long. The algorithm never enumerates it: the finite part enters only through the commutant
$\mathcal E$ and the minimality argument of Lemma~\ref{lem:abext}, which replaces averaging over the finite group.
The finiteness test of Lemma~\ref{lem:finite}(1) is applied only to the given generators; the bounds of
Lemma~\ref{lem:finite}(2) and the Schreier generators are used only in the analysis
(Lemmas~\ref{lem:apriori} and~\ref{lem:Upp}), not in the algorithm.

(ii) Unipotent parts of short words do not suffice to generate the Lie algebra of $U$. Let
$\Gamma\le\GL_{n+2}(\Q)$ be the affine Weyl group of type $\tilde A_n$, acting on $\Q^{n+1}$ by the
affine maps $s_i$ (transposition of coordinates $i,i+1$, $1\le i\le n$) and
$s_0:x\mapsto(1\ n{+}1)x+e_1-e_{n+1}$. Any word of length at most $n$ in $s_0,\ldots,s_n$ involves at
most $n$ of the $n+1$ Coxeter generators, hence lies in a proper parabolic subgroup, which is finite; so it
has finite order and trivial unipotent part. Yet $\Gamma$ is quasi-unipotent and virtually abelian with
$U$ the $n$-dimensional group of translations by vectors with coordinate sum zero, and
$|\cG/\cG^\circ|=(n+1)!$. (Coxeter elements, of length $n+1$, do have nontrivial unipotent part, so the
example only shows that the necessary word length grows with $d$; we know no polynomial bound in
general, and the algorithm avoids the question by working with Levi conjugators rather than with group
elements. Replacing translations by $\operatorname{diag}(2,1,\ldots,1,\tfrac12)$ gives the same phenomenon
for the torus.)

(iii) Testing (a) is at least as hard as testing finiteness: for $\Gamma$ inside the orthogonal group
of the anisotropic form $x_1^2+x_2^2+x_3^2$ over $\Q$, (a) holds iff $\Gamma$ is finite, since the identity
component of the closure is trivial, a torus without nontrivial rational characters (a split torus would
have rational isotropic weight vectors), or $SO_3$.

(iv) Beals~\cite{Bea99} decides virtual solvability of finitely generated linear groups over number fields in
polynomial time by the same centralizer construction as in Lemma~\ref{lem:atest} (his Lemma~3.1 and
Section~5), and states that virtual abelianness is decidable in polynomial time; we have not relied on the
latter statement, whose proof there is only sketched. Lemma~\ref{lem:atest} adds the splitting test, which
is what distinguishes CCRA-computability: the series with $M=\left(\begin{smallmatrix}1&1\\1&2\end{smallmatrix}\right)$
has a virtually abelian group but a non-split torus, so it is not CCRA-computable.
\end{remark}

\subsection{Polynomial heights: a deterministic polynomial-time algorithm}\label{ssec:detpoly}

We now prove that $B_f$ is bounded by a polynomial in the input length. Together with Theorem~\ref{thm:polytime}(2)
this shows that the algorithm of Section~\ref{ssec:layers}, unchanged, runs in deterministic
polynomial time on all invertible series. Only the analysis is new.

\begin{theorem}\label{thm:detpoly}
There is a polynomial $p$ such that, for every invertible series $f$ satisfying~(a), the Lie algebra
$\fu$ of $U=R_u(\cG)$ is spanned by matrices whose entries, in the adapted basis of step~1, have bit
length at most $p(d,k,\sigma)$. Consequently $\operatorname{ht}(\fu_j)\le p'(d,k,\sigma)$ for all
$j$, for a polynomial $p'$; in particular $B_f\le p'(d,k,\sigma)$, and CCRA-computability of
invertible series is decidable in deterministic polynomial time.
\end{theorem}

\paragraph{Conventions.} Throughout, (a) holds and $S:=H^\circ$, a $\Q$-split torus contained in
$L_P$. A rational matrix or vector has \emph{polynomial size} if the bit lengths of its entries are
bounded by a fixed polynomial in $d,k,\sigma$ (independent of $f$); the polynomial may change from
line to line, finitely many times. We use three standard facts. (F1) A subspace of $\Q^N$ spanned by
vectors of bit length $\le\beta$ has height at most $\operatorname{poly}(N)\,(\beta+1)$ (Cramer's rule
and Hadamard's bound). (F2) If $V_0\subseteq\Q^N$ is spanned by vectors of bit length $\le\beta$ and
$A_1,\ldots,A_s$ are $N\times N$ matrices with entries of bit length $\le\beta$, then the smallest
subspace containing $V_0$ and stable under all $A_i$ is spanned by the vectors $A_{i_1}\cdots
A_{i_r}v$ with $r<N$ and $v$ from the spanning set; a product of $r$ such matrices has a common
denominator of at most $rN^2\beta$ bits, so these vectors have bit length $O(N^3(\beta+\log N))$.
(F3) If $V\subseteq\fu_P$ is spanned by matrices of bit length $\le\beta$, the Lie algebra generated
by $V$ is spanned by the left-normed brackets of fewer than $m$ of them (brackets of length $m$
vanish), which have bit length $O(md^2(\beta+\log d))$; if $V$ is stable under a family of Lie algebra
automorphisms or derivations of $\fu_P$, so is the Lie algebra it generates.

\paragraph{Weight basis.} The algebra $C=\Q[H^\circ]=\Q[S]$ is computed in Lemma~\ref{lem:atest} (by
the proof of that lemma, under~(a) the computed algebra equals $\Q[H^\circ]$, which lies in the span of
$H\subseteq L_P$); it
is commutative, split and block diagonal, and its primitive idempotents $\varepsilon_\chi$ (one for
each weight $\chi$ of $S$ on $\Q^d$) have polynomial size: $C$ has a basis of polynomial size (it is a
centralizer computed by linear algebra), a combination of the basis elements with coefficients in
$\{0,\ldots,d^2\}$ separates the at most $d$ weights (by the box argument used in Lemma~\ref{lem:v0}), and each
$\varepsilon_\chi$ is a Lagrange interpolation polynomial of degree $<d$ in that element, whose
rational eigenvalues are found by rational root finding. Refining the adapted basis on each $E_p$ by bases of the
$\varepsilon_\chi E_p$ changes the basis by a block diagonal matrix of polynomial size; we work in
this refined basis and transform back at the end, which preserves polynomial size of spanning sets.
Let $X_S$ be the set of weights, $V_\chi=\varepsilon_\chi\Q^d$ and $n_\chi=\dim V_\chi$; each basis
vector lies in some $V_\chi$, and $S$ is contained in the diagonal torus
$D=\{\sum_\chi x_\chi\varepsilon_\chi\}\cong\mathbb G_m^{X_S}$. Write $\mathfrak z\subseteq\fu_P$ for
the span of the elementary matrices $E_{rc}\in\fu_P$ with $r,c$ in the same weight space, and
$\fu_P^{\ne0}$ for the span of the other elementary matrices in $\fu_P$. These are the sums of the
zero and of the nonzero weight spaces of $\operatorname{Ad}(S)$ on $\fu_P$; both are graded by distance,
$\mathfrak z$ is an associative subalgebra, $Z:=C_{U_P}(S)=I+\mathfrak z=\exp\mathfrak z$, and
$\operatorname{Ad}(h)$ preserves $\mathfrak z$ and $\fu_P^{\ne0}$ for $h\in H$, because $H$ normalizes $S$ and
hence permutes the $V_\chi$. For $h\in H$ write $\pi_h$ for the induced permutation of $X_S$,
$hV_\chi=V_{\pi_h\chi}$. Every element of $H$ is semisimple, since a power lies in the torus $S$.

An element $Y=\sum_\chi y_\chi\varepsilon_\chi\in\operatorname{Lie}(S)$ is \emph{regular} if
$y_\chi\ne y_\psi$ whenever $\chi\ne\psi$ and $\fu_P$ contains an elementary matrix $E_{rc}$ with
$r\in V_\chi$, $c\in V_\psi$. Then $\operatorname{ad}(Y)$ acts on $E_{rc}$ by $y_\chi-y_\psi$, so it is invertible on
$\fu_P^{\ne0}$ and zero on $\mathfrak z$, and a matrix $v\in U_P$ commutes with $Y$ iff $v\in Z$, that is,
\begin{equation}\label{eq:regcent}
C_{U_P}(Y)=Z\qquad\text{for regular }Y.
\end{equation}
If $Y$ is regular, so is $\operatorname{Ad}(h)Y$ for $h\in H$. Since $\chi-\psi$ is a nontrivial character of
$S$ for $\chi\ne\psi$, its differential is a nonzero linear form on $\operatorname{Lie}(S)$; so the regular elements
form the complement of at most $d^2$ hyperplanes of $\operatorname{Lie}(S)$.

\begin{lemma}[Tame iterations]\label{lem:iter}
Let $\fu_P^{(t)}$ denote the matrices supported at distances $\ge t$.
\begin{enumerate}
\item Let $h\in L_P(\Q)$ be semisimple and $g\in hU_P(\Q)$, both of polynomial size. There is
$u\in U_P(\Q)$ of polynomial size with $u^{-1}gu\in h\,C_{U_P}(h)$. Consequently $g_s=uhu^{-1}$ and the
closure of $\langle g_s\rangle$ is $u\,\overline{\langle h\rangle}\,u^{-1}$.
\item Let $Y\in\operatorname{Lie}(S)$ be regular and $X\in Y+\fu_P$, both of polynomial size. There is
$u\in U_P(\Q)$ of polynomial size with $\operatorname{Ad}(u^{-1})X\in Y+\mathfrak z$. Consequently
$X_s=\operatorname{Ad}(u)Y$.
\item The map $\mu:\exp(\fu_P^{\ne0})\times Z\to U_P$, $(e^y,z)\mapsto e^yz$, is an isomorphism of
varieties defined over $\Q$ (here $\exp(\fu_P^{\ne0})$ is a closed subvariety of $U_P$, in general not a
subgroup), and if $w\in U_P(\Q)$ has polynomial size, so does the $y$ with
$w\in e^yZ$.
\end{enumerate}
\end{lemma}
\begin{proof}
All three are proved by an iteration over $t=1,\ldots,m-1$ that, at step $t$, removes an unwanted
distance-$t$ component without changing components at smaller distances. We use repeatedly that for
$A\in\fu_P^{(t)}$ and $B\in\fu_P$ every Lie bracket involving $A$ and at least one further factor
from $\fu_P$ lies in $\fu_P^{(t+1)}$, and the Baker--Campbell--Hausdorff formula, which is finite on
$\fu_P$.

(1) Since $h$ is semisimple and preserves each $\mathfrak v_t$, we have
$\mathfrak v_t=\ker(1-\operatorname{Ad}(h^{-1}))\oplus\operatorname{im}(1-\operatorname{Ad}(h^{-1}))$ on $\mathfrak v_t$. Start with $u=I$. At
step $t$ let $R=\log(h^{-1}u^{-1}gu)\in\fu_P$ and suppose that its components at distances $<t$ lie in
the kernels. Let $Q_t$ be the component of $R_t$ in the image, choose $Z_t\in\mathfrak v_t$ with
$(1-\operatorname{Ad}(h^{-1}))Z_t=-Q_t$, and replace $u$ by $ue^{Z_t}$. The new element is
$h^{-1}e^{-Z_t}u^{-1}gue^{Z_t}=e^{-\operatorname{Ad}(h^{-1})Z_t}e^{R}e^{Z_t}$, whose logarithm is
$R+(1-\operatorname{Ad}(h^{-1}))Z_t$ modulo $\fu_P^{(t+1)}$. Thus the components at distances $<t$ are
unchanged and the distance-$t$ component now lies in the kernel. After step $m-1$ we have
$u^{-1}gu=he^R$ with $\operatorname{Ad}(h)R=R$, so $h$ and $e^R$ commute; this is the Jordan decomposition of
$u^{-1}gu$, whence $g_s=uhu^{-1}$.

(2) Start with $u=I$. At step $t$ let $R=\operatorname{Ad}(u^{-1})X-Y\in\fu_P$, suppose that its components at
distances $<t$ lie in $\mathfrak z$, let $Q_t$ be the component of $R_t$ in $\fu_P^{\ne0}$, let
$Z_t\in\mathfrak v_t\cap\fu_P^{\ne0}$ be the solution of $[Z_t,Y]=Q_t$ (entrywise division by the nonzero
numbers $y_\psi-y_\chi$), and replace $u$ by $ue^{Z_t}$. Then
$\operatorname{Ad}(e^{-Z_t})(Y+R)=Y+R-[Z_t,Y]-[Z_t,R]+\tfrac12[Z_t,[Z_t,Y+R]]-\cdots$, which is
$Y+R-Q_t$ modulo $\fu_P^{(t+1)}$. After step $m-1$, $\operatorname{Ad}(u^{-1})X=Y+N$ with $N\in\mathfrak z$
nilpotent and $[Y,N]=0$, so $X_s=\operatorname{Ad}(u)Y$.

(3) Start with $y=0$. At step $t$ let $R=\log(e^{-y}w)$, suppose that its components at distances
$<t$ lie in $\mathfrak z$, and replace $y$ by $y+Q_t$, where $Q_t$ is the component of $R_t$ in
$\fu_P^{\ne0}$. Since $\log(e^{-y-Q_t}e^{y})\equiv-Q_t$ modulo $\fu_P^{(t+1)}$, the new logarithm is
$R-Q_t$ modulo $\fu_P^{(t+1)}$. At the end $\log(e^{-y}w)\in\mathfrak z$, i.e.\ $e^{-y}w\in Z$; so $\mu$
is onto, with an inverse given by polynomials. It is injective: if $e^{y}z=e^{y'}z'$ with $y\ne y'$,
let $t$ be the least distance with $y_t\ne y'_t$; then $\log(e^{-y'}e^{y})\equiv y-y'$ modulo
$\fu_P^{(t+1)}$ has the nonzero distance-$t$ component $y_t-y'_t\in\fu_P^{\ne0}$, but it lies in
$\log Z=\mathfrak z$.

\emph{Sizes.} We use the tame calculus of the proof of Lemma~\ref{lem:sizes}. Let $\delta$ be the
product of $m!$, of $2^{\beta}$ where $\beta$ bounds the bit lengths of the input data ($h^{\pm1}$ and
$h^{-1}g$ in (1), $Y$ and $X-Y$ in (2), $w$ in (3)), of their denominators, and of the numerators and
denominators of the following fixed linear maps of $\mathfrak v_t$ ($1\le t<m$): in (1) the map $A_t$ that
projects onto $\operatorname{im}(1-\operatorname{Ad}(h^{-1}))$ along the kernel and then inverts $1-\operatorname{Ad}(h^{-1})$ there; in (2) the
entrywise division by the numbers $y_\psi-y_\chi$; in (1)--(3) the projections onto $\fu_P^{\ne0}$ and
onto a distance component. These maps have entries of polynomial bit length (for $A_t$ by Cramer's
rule applied to the matrix of $\operatorname{Ad}(h^{-1})$ on $\mathfrak v_t$), so $\log\delta$ is polynomial. In
each step the current $u$ (resp.\ $y$) is updated incrementally ($u\mapsto ue^{Z_t}$, $y\mapsto y+Q_t$),
and the quantities $R$, $Q_t$, $Z_t$, $e^{\pm Z_t}$ and the updated $u$ (in~(1),
$h^{-1}u^{-1}gu=(h^{-1}u^{-1}h)(h^{-1}g)u$; in~(2),
$\operatorname{Ad}(u^{-1})X-Y=(u^{-1}-1)Y+Y(u-1)+(u^{-1}-1)Y(u-1)+u^{-1}(X-Y)u$) are obtained from the previous
ones and the input by a bounded number of the operations of that calculus: sums and products,
$\exp(X)-1$, $\log(1+X)$, $(1+X)^{-1}-1$, conjugation by $h^{\pm1}$, products with $Y$, projections, and
application of one of the fixed maps above to a single distance component. As in the proof of
Lemma~\ref{lem:sizes} (where the same structure occurs: each new quantity is obtained from the
previous iterate and the fixed input by $O(1)$ operations, so every quantity is the result of a chain
of $O(m)$ operations, and two-argument operations are applied after raising both arguments to the
larger $(c,\kappa)$ by monotonicity), each such operation maps $(c,\kappa)$-tame matrices to
$(c+2,F\kappa)$-tame ones with $\log F$ polynomial (for a fixed map with entries of absolute value
$\le\delta$ and denominators dividing $\delta$, applied to a distance-$t$ component, one may take
$F=\dim(\mathfrak v_t)\,\delta$). Starting from $(2,d\delta^2)$-tameness of the input, after the $O(m)$
operations of the whole iteration all matrices are $(O(m),F^{O(m)}d\delta^2)$-tame, and an entry at
distance $s\le m$ of a $(c,\kappa)$-tame matrix has bit length at most $2cs\log\delta+s\log\kappa+2$,
which is polynomial. The block diagonal parts ($h$, $Y$, and $u$ having identity diagonal blocks) are
fixed throughout.
\end{proof}

\paragraph{The torus from short words.} For a word $w$ in the letters $M_a^{\pm1}$ write $M_w$ for its
matrix, $|w|$ for its length, and $S_w$ for the identity component of the closure of
$\langle b(M_w)\rangle$; it is a subtorus of $S$, since $b(M_w)^e\in S$ for $e=|H/S|$.

For a subgroup $\Delta$ of $D(\Q)$ let $\Lambda(\Delta)\subseteq\Q^{X_S}$ be the span of the vectors
$(v_p(t_\chi))_{\chi\in X_S}$ over all $t=\sum_\chi t_\chi\varepsilon_\chi\in\Delta$ and all primes $p$,
where $v_p$ is the $p$-adic valuation. The identity component of the closure of $\Delta$ is cut out in
$D$ by the characters $x\mapsto\prod_\chi x_\chi^{a_\chi}$ ($a\in\Z^{X_S}$) that take root-of-unity
values, i.e.\ the values $\pm1$, on $\Delta$~\cite[8.2, 8.12]{Bor91}, and $\prod_\chi t_\chi^{a_\chi}=\pm1$
iff $\sum_\chi a_\chi v_p(t_\chi)=0$ for all $p$. Hence, identifying $y\in\Q^{X_S}$ with
$\sum_\chi y_\chi\varepsilon_\chi$,
\begin{equation}\label{eq:lietorus}
\operatorname{Lie}\bigl(\overline\Delta^{\,\circ}\bigr)=\Lambda(\Delta).
\end{equation}

\begin{lemma}[Short words generate the torus]\label{lem:torus}
$\operatorname{Lie}(S)=\sum_w\operatorname{Lie}(S_w)$, the sum over the words of length at most $2d^2+2d$. For each such
word, $\operatorname{Lie}(S_w)$ is spanned by vectors of polynomial size.
\end{lemma}
\begin{proof}
For $g\in b(\Gamma)$ and $\chi\in X_S$ let $g_\chi:V_{\pi_g^{-1}\chi}\to V_\chi$ be the restriction
of $g$, a square matrix in the weight basis. Let $\mathcal P$ be the finite set of primes dividing a
numerator or denominator of some $\det((b_a^{\pm1})_\chi)$. For $g\in b(\Gamma)$ define
$c(g)\in\mathcal Q:=\Q^{X_S}\otimes\Q^{\mathcal P}$ by $c(g)_{\chi,p}=v_p(\det g_\chi)$. Since
$(gh)_\chi=g_\chi h_{\pi_g^{-1}\chi}$, we have the cocycle identity $c(gh)=c(g)+\pi_gc(h)$, where
$(\pi_gu)_\chi=u_{\pi_g^{-1}\chi}$; in particular every $\det g_\chi$ is a product of the generator
determinants, so its valuations outside $\mathcal P$ vanish, and $|c(b(M_w))_{\chi,p}|\le|w|\beta_0$, where
$\beta_0$ bounds the bit lengths of the determinants $\det((b_a^{\pm1})_\chi)$ in the weight basis; it is
polynomial in $d$ and $\sigma$, by Hadamard's bound. For
$t\in S$ we have $\det t_\chi=t_\chi^{n_\chi}$, so with $\mathsf N=\operatorname{diag}(n_\chi)$,
\begin{equation}\label{eq:cval}
\bigl(v_p(t_\chi)\bigr)_\chi=\mathsf N^{-1}c(t)_{\cdot,p}\qquad(t\in S\cap b(\Gamma)).
\end{equation}
Let $K_\pi=\{g\in\Gamma:\pi_{b(g)}=1\}$, a normal subgroup of finite index on which $g\mapsto c(b(g))$ is
a homomorphism, and $\mathsf W=\spn\{c(b(g)):g\in K_\pi\}$. For $g\in K_\pi$ some power $g^e$ has $b(g^e)\in S$ and
$c(b(g^e))=e\,c(b(g))$; conversely $S\cap b(\Gamma)\subseteq b(K_\pi)$. For $u\in\mathcal Q$ let
$\operatorname{col}(u)\subseteq\Q^{X_S}$ be the span of the vectors $u_{\cdot,p}$ ($p\in\mathcal P$), and for a
subspace $\mathsf W_1$ let $\operatorname{col}(\mathsf W_1)=\sum_{u\in\mathsf W_1}\operatorname{col}(u)$. Since $S\cap b(\Gamma)$ is dense in $S$
(Lemma~\ref{lem:monoid}), \eqref{eq:lietorus} and \eqref{eq:cval} give
$\operatorname{Lie}(S)=\operatorname{col}(\mathsf N^{-1}\mathsf W)$. Similarly, for a word $w$ let $n_w$ be the order of
$\pi_{b(M_w)}$ and $\nu_w:=c(b(M_w)^{n_w})=\sum_{i<n_w}\pi_{b(M_w)}^ic(b(M_w))$; as
$S_w$ is the identity component of the closure of $\langle b(M_w)^{n_we}\rangle$ and
$c(b(M_w)^{n_we})=e\nu_w$, we get $\operatorname{Lie}(S_w)=\operatorname{col}(\mathsf N^{-1}\nu_w)$. Its spanning vectors
have entries of absolute value at most $n_w|w|\beta_0$ with $n_w\le d!$, hence polynomial size. Moreover, for
words $x,w$, using $c(g^{-1})=-\pi_g^{-1}c(g)$ and $\pi_{b(M_w)^{n_w}}=1$,
\begin{equation}\label{eq:conjnu}
\nu_{xwx^{-1}}=c\bigl(b(M_x)\,b(M_w)^{n_w}\,b(M_x)^{-1}\bigr)=\pi_{b(M_x)}\nu_w .
\end{equation}

Let $\mathsf W'$ be the span of the $\pi_{b(M_x)}\nu_w$ with $|w|\le2d-1$ and $x$ a positive word with $|x|\le d^2$. Then
$\mathsf W'\subseteq\mathsf W$, because $M_w^{n_w}\in K_\pi$. For a fixed $u$, let $O_i$ be the span of the
$\pi_{b(M_x)}u$ with $x$ a positive word of length $\le i$ in the letters $M_a$. Then $O_i\subseteq O_{i+1}$,
and $O_i=O_{i+1}$ implies that $O_i$ is stable under all $\pi_{b_a}$, hence under the finite group they
generate, which contains all $\pi_g$ ($g\in b(\Gamma)$). Since $O_i$ lies in the image of the linear map
$\spn\{\pi_g\}\to\mathcal Q$, $A\mapsto Au$, its dimension is at most $\dim\spn\{\pi_g\}\le|X_S|^2\le d^2$;
so $O_{d^2}$ is the span of all $\pi_gu$, and $\mathsf W'$ is stable under all $\pi_g$. We claim
$\mathsf W'=\mathsf W$. Otherwise let $\mathcal Q_1$ be a $\pi$-stable complement of $\mathsf W'$ in $\mathcal Q$
(the $\pi_g$ form a finite group) and $q$ the equivariant projection onto $\mathcal Q_1$ along $\mathsf W'$.
Then $\bar c=q\circ c$ is again a cocycle with values in the permutation module $\mathcal Q$, and
$q(\mathsf W)\ne0$. Unwinding the cocycle identity,
for $g_i=b(M_{a_i}^{\eps_i})$,
\[
\bar c(g_1\cdots g_s)_\chi=\sum_{i=1}^s\bar c(g_i)_{\chi_{i-1}},\qquad\chi_0=\chi,\quad
\chi_i=\pi_{g_i}^{-1}\chi_{i-1}.
\]
Let $\Omega$ be the graph with vertex set $X_S$ and, for each letter $a$ and each $\chi$, an edge from
$\chi$ to $\pi_{b_a}^{-1}\chi$ labelled $\bar c(b_a)_\chi\in\Q^{\mathcal P}$. Since
$\bar c(b_a^{-1})_\chi=-\bar c(b_a)_{\pi_{b_a}\chi}$, the letter $a^{-1}$ traverses an edge backwards
with the opposite label. Thus words correspond to walks, $\bar c(b(M_w))_\chi$ is the label of the walk
of $w$ from $\chi$, and $w\in K_\pi$ iff all these walks are closed. As $q(\mathsf W)\ne0$, some $g\in K_\pi$
and some $\chi$ give a closed walk with nonzero label. The label of a walk is the sum of the labels of the
edges it traverses, counted with sign, so it depends only on the vector in $\Z^{E(\Omega)}$ of signed
traversal counts. For a closed walk this vector lies in the cycle space of $\Omega$, which has a
$\Z$-basis given by the fundamental cycles of a spanning forest (the coefficient of a fundamental cycle
is the signed count of its non-forest edge). Hence some fundamental cycle has nonzero label. Read from the root $r$ of its tree, it is a closed walk of length at most $2(|X_S|-1)+1\le2d-1$,
i.e.\ the walk from $r$ of a word $w$ with $|w|\le2d-1$ and $\pi_{b(M_w)}r=r$. Then
$q(\nu_w)_r=\sum_{i<n_w}\bar c(b(M_w))_{\pi^{-i}_{b(M_w)}r}=n_w\,\bar c(b(M_w))_r\ne0$, so
$\nu_w\notin\mathsf W'$, a contradiction.

Hence $\operatorname{Lie}(S)=\operatorname{col}(\mathsf N^{-1}\mathsf W')=\sum_{x,w}\operatorname{col}(\mathsf N^{-1}\nu_{xwx^{-1}})
=\sum_{x,w}\operatorname{Lie}(S_{xwx^{-1}})$ by~\eqref{eq:conjnu}, where $|xwx^{-1}|\le2d^2+2d-1$.
\end{proof}

\begin{lemma}[A maximal torus of polynomial size]\label{lem:v0}
There are a regular $Y\in\operatorname{Lie}(S)$ and $v_0\in U_P(\Q)$, both of polynomial size, such that
$v_0Sv_0^{-1}\subseteq\cG^\circ$.
\end{lemma}
\begin{proof}
Let $w$ be a word with $|w|\le2d^2+2d$. Lemma~\ref{lem:iter}(1), applied to $h=b(M_w)$ and $g=M_w$,
gives $u_w\in U_P(\Q)$ of polynomial size with $(M_w)_s=u_wb(M_w)u_w^{-1}$. Hence the torus
$\Theta_w$, the identity component of the closure of $\langle(M_w)_s\rangle$, equals $u_wS_wu_w^{-1}$;
it lies in $\cG^\circ$, because $(M_w)_s\in\cG$ (A2) and $\Theta_w$ is connected. So
$\operatorname{Ad}(u_w)\operatorname{Lie}(S_w)\subseteq\operatorname{Lie}(\cG^\circ)$. Let $Y_{w,1},Y_{w,2},\ldots$ be spanning vectors of
$\operatorname{Lie}(S_w)$ of polynomial size (Lemma~\ref{lem:torus}). Altogether they span $\operatorname{Lie}(S)$, so we can choose
among them a basis $Y_1,\ldots,Y_n$ of $\operatorname{Lie}(S)$, with $n=\dim S\le d$, where $Y_i$ is a spanning vector of
$\operatorname{Lie}(S_{w_i})$; only these $n$ words are used below. For each of the at most $d^2$ pairs $(\chi,\psi)$ in the definition of regularity,
$c\mapsto(y_\chi-y_\psi)\bigl(\sum_i c_iY_i\bigr)$ is a nonzero linear form on $\Q^n$, which
vanishes on at most $(d^2+1)^{n-1}$ points of the box $\{0,\ldots,d^2\}^n$. So some $c$ in the box
avoids all of them. Put $Y=\sum_ic_iY_i$, which is regular, and
$X=\sum_ic_i\operatorname{Ad}(u_{w_i})Y_i\in\operatorname{Lie}(\cG^\circ)$. Both have polynomial size, being sums of $n\le d$ terms of
polynomial size, and $b(X)=Y$, because
$\operatorname{Ad}(u)B-B\in\fu_P$ for $u\in U_P$ and block diagonal $B$. Lemma~\ref{lem:iter}(2) gives
$v_0\in U_P(\Q)$ of polynomial size with $X_s=\operatorname{Ad}(v_0)Y$.

It remains to show $v_0Sv_0^{-1}\subseteq\cG^\circ$. Let $T'$ be a maximal torus of $\cG^\circ$. Then
$b(T')=b(\cG^\circ)=H^\circ=S$ (A2), since $\cG^\circ=T'U$ and $b(U)=1$. The tori $T'$ and $S$ are maximal
tori of the connected solvable group $b^{-1}(S)=S\ltimes U_P$, so $T'=vSv^{-1}$ for some $v\in U_P$
\cite[10.6]{Bor91}. Put $\fu'=\operatorname{Ad}(v^{-1})\fu$ and $X'=\operatorname{Ad}(v^{-1})X$. As
$\operatorname{Lie}(\cG^\circ)=\operatorname{Lie}(T')\oplus\fu$, we have $X'\in\operatorname{Lie}(S)\oplus\fu'$, and comparing images under $b$
gives $X'=Y+N$ with $N\in\fu'$. Now $\fu'$ is $\operatorname{Ad}(S)$-stable, so it is the sum of its intersections
with $\mathfrak z$ and with $\fu_P^{\ne0}$, and $\operatorname{ad}(Y)$ maps $\fu'\cap\fu_P^{\ne0}$ bijectively onto itself.
If the component $N^{\ne0}$ of $N$ in $\fu_P^{\ne0}$ lies in $\fu_P^{(t)}$, let $\Xi\in\fu'\cap\fu_P^{\ne0}$
with $[\Xi,Y]=N^{\ne0}$; then $\Xi\in\fu_P^{(t)}$, since $\operatorname{ad}(Y)$ preserves distances, and
$\operatorname{Ad}(e^{-\Xi})(Y+N)=Y+N-N^{\ne0}+R$ with $R=-[\Xi,N]+\frac12[\Xi,[\Xi,Y+N]]-\cdots$, which lies in
$\fu'\cap\fu_P^{(t+1)}$ (every term is a bracket of $\Xi\in\fu_P^{(t)}$ with an element of $\fu'$, using
$[\Xi,Y]=N^{\ne0}\in\fu'$). Replacing $N$ by $N-N^{\ne0}+R$, the new component in $\fu_P^{\ne0}$ lies in
$\fu_P^{(t+1)}$. After fewer than $m$ such steps we obtain $\upsilon=e^{\Xi_1}\cdots e^{\Xi_r}\in\exp(\fu')=v^{-1}Uv$ with
$\operatorname{Ad}(\upsilon^{-1})X'\in Y+\mathfrak z$, so $X'_s=\operatorname{Ad}(\upsilon)Y$ as in Lemma~\ref{lem:iter}(2). Thus
$v''=v\upsilon\in U_P$ satisfies $\operatorname{Ad}(v'')Y=X_s$, and
$v''Sv''^{-1}=(v\upsilon v^{-1})T'(v\upsilon v^{-1})^{-1}\subseteq\cG^\circ$, as $v\upsilon v^{-1}\in U$. Finally
$\operatorname{Ad}(v_0)Y=\operatorname{Ad}(v'')Y$, so $v''^{-1}v_0\in C_{U_P}(Y)=Z$ by~\eqref{eq:regcent}, and $Z$ centralizes $S$;
hence $v_0Sv_0^{-1}=v''Sv''^{-1}\subseteq\cG^\circ$.
\end{proof}

From now on we conjugate by $v_0$: we replace $M_a$ by $v_0^{-1}M_av_0$, and $\cG$, $U$, $\fu$ by their
conjugates. This changes neither the $b_a$ nor polynomial size, and now $S\subseteq\cG^\circ$.

\begin{lemma}[Moving the generators into the normalizer of $S$]\label{lem:xhat}
For each $a$ there is $\hat x_a\in U(\Q)$ of polynomial size such that $m_a:=\hat x_a^{-1}M_a$ normalizes
$S$. Moreover $m_a\in\cG$, $m_a$ has polynomial size, and $b(m_a)=b_a$.
\end{lemma}
\begin{proof}
As $S\subseteq\cG$ normalizes $U$, the Lie algebra $\fu$ is $\operatorname{Ad}(S)$-stable, so
$\fu=\fu^0\oplus\fu^{\ne0}$ with $\fu^0=\fu\cap\mathfrak z$ and $\fu^{\ne0}=\fu\cap\fu_P^{\ne0}$, and
$C_U(S)=\exp(\fu^0)$, since $e^X$ commutes with $s\in S$ iff $\operatorname{Ad}(s)X=X$. We claim
$U=\mu\bigl(\exp(\fu^{\ne0})\times C_U(S)\bigr)$. The set $\mu^{-1}(U)$ is closed and irreducible of dimension
$\dim U$, because $\mu$ is an isomorphism (Lemma~\ref{lem:iter}(3)). It contains
$\exp(\fu^{\ne0})\times\exp(\fu^0)$, which is closed and irreducible of dimension
$\dim\fu^{\ne0}+\dim\fu^0=\dim U$; so the two sets are equal.

The torus $S$ is a maximal torus of $\cG^\circ$, since the maximal tori of $\cG^\circ$ map isomorphically onto
$b(\cG^\circ)=S$. As $\cG^\circ$ is normal in $\cG$, $M_aSM_a^{-1}$ is a maximal torus of $\cG^\circ$ too, so
$M_aSM_a^{-1}=xSx^{-1}$ for some $x\in U$ \cite[10.6]{Bor91}. Write $x=e^yz$ with $y\in\fu^{\ne0}$ and
$z\in C_U(S)$, and put $\hat x_a=e^y\in U$. Then $\hat x_aS\hat x_a^{-1}=M_aSM_a^{-1}$, so
$m_a=\hat x_a^{-1}M_a\in\cG$ normalizes $S$. Let $Y$ be as in Lemma~\ref{lem:v0} and $Y_a=\operatorname{Ad}(b_a)Y$,
which is regular. The element $\operatorname{Ad}(m_a)Y$ lies in $\operatorname{Lie}(S)$, and elements of $\operatorname{Lie}(S)$ are determined by
their images under $b$; so $\operatorname{Ad}(m_a)Y=Y_a$, i.e.\ $\operatorname{Ad}(\hat x_a)Y_a=\operatorname{Ad}(M_a)Y$. By~\eqref{eq:regcent}
and the injectivity of $\mu$, $\hat x_a$ is the only element of $\exp(\fu_P^{\ne0})$ with this property.
The linear system $(I+N)Y_a=\operatorname{Ad}(M_a)Y\,(I+N)$ in $N\in\fu_P$ has rational coefficients of polynomial size
and the solution $\hat x_a-I$, a priori over $\Qb$; so it is consistent over $\Q$. Let $w_a$ be $I$ plus its
reduced echelon solution (free variables set to zero), which is rational of polynomial size by Cramer's
rule and Hadamard's bound. Then $w_a\in\hat x_aC_{U_P}(Y_a)=\hat x_aZ$ by~\eqref{eq:regcent}, so $\hat x_a=e^y$ with
$y$ the component of $w_a$ given by Lemma~\ref{lem:iter}(3). Hence $\hat x_a$ is rational of polynomial
size, and so is $m_a$. Finally $b(m_a)=b(M_a)=b_a$, because $\hat x_a\in U_P$.
\end{proof}

\begin{lemma}[The part of $U$ centralizing $S$]\label{lem:Upp}
Let $N'$ be the closure of the group generated by $S$ and $m_1,\ldots,m_k$, and $U''=N'\cap U_P$.
Then $U''\subseteq U$, and $\operatorname{Lie}(U'')$ is spanned by matrices of polynomial size.
\end{lemma}
\begin{proof}
$N'\subseteq\cG$ normalizes $S$, and $U''\subseteq\cG\cap U_P=U$. An element $u\in U_P$ normalizing $S$
centralizes it, because $usu^{-1}\in S$ and $b(usu^{-1})=s$; so $U_P\cap N_{\GL_d}(S)=Z$.

The algebra $C(S)=\bigoplus_\chi\operatorname{End}(V_\chi)$ (the centralizer of $S$ in $M_d(\Q)$) is stable under
conjugation by $b(\Gamma)$, and $S$ acts trivially on it. Let $\Psi$ be the group of these conjugation
operators. It is finite, as its kernel contains $S\cap b(\Gamma)$, which has finite index in $b(\Gamma)$.
By Lemma~\ref{lem:finite}(2), applied to the generators $\operatorname{Ad}(b_a)|_{C(S)}$ in dimension $\dim C(S)=\sum_\chi n_\chi^2\le d^2$,
$|\Psi|\le(2d^2+1)^{d^4}$, and the elements of $\Psi$ have entries of polynomial bit length with a common
denominator of polynomial bit length. The centralizer of $C(S)$ in $\GL_d$ is $D$, so $\operatorname{Ad}(b(g))$ is
trivial on $C(S)$ iff $b(g)\in D$.

Let $Y_N=\{g\in N':b(g)\in D\}$ and $\varphi(g)=b(g)^{-1}g$ for $g\in Y_N$. For $g\in Y_N$ the element
$n=b(g)^{-1}g\in U_P$ normalizes $S$, since $g$ and $b(g)$ do; so $n\in Z\subseteq C(S)$, and $n$ commutes
with $b(g)$, which is central in $C(S)$. As $b(g)$ is semisimple and $n$ unipotent, $g_s=b(g)$ and $g_u=n$,
and both lie in $N'$ (A2). Hence $\varphi(g)\in N'\cap U_P=U''$. Moreover $\varphi$ is a homomorphism:
$\varphi(gh)=b(h)^{-1}\varphi(g)b(h)\varphi(h)=\varphi(g)\varphi(h)$, as $\varphi(g)\in Z\subseteq C(S)$ commutes with
$b(h)\in D$. Since $U''\subseteq Y_N$ and $\varphi$ is the identity on $U''$, $U''=\varphi(Y_N)$.

Let $\Gamma''=\langle m_1,\ldots,m_k\rangle$, let $N'_0$ be its closure, and let $K''=\Gamma''\cap Y_N$, the kernel of
$\Gamma''\to\Psi$, $g\mapsto\operatorname{Ad}(b(g))|_{C(S)}$; it has index at most $|\Psi|$. The product $N'_0S$ is a closed
subgroup containing $\Gamma''$ and $S$, so $N'=N'_0S$, and $Y_N=(N'_0\cap Y_N)S$ as $S\subseteq Y_N$. Since
$\varphi(S)=1$, $U''=\varphi(N'_0\cap Y_N)$. The closure of $K''$ is a closed subgroup of finite index in
$N'_0\cap Y_N$, so its image under $\varphi$ is a closed subgroup of finite index of the connected group $U''$,
hence equal to $U''$. Thus $U''$ is the closure of the group generated by $\varphi(K'')$, and so by the
$\varphi(\varsigma)$ for the Schreier generators $\varsigma=r_\psi m_ar_{\psi'}^{-1}$ of $K''$, where $r_\psi$ is a
shortest positive word in the $m_a$ mapping to $\psi\in\Psi$. These are words of length $s<2|\Psi|$ in the
$m_a^{\pm1}$. As in the proof of Lemma~\ref{lem:apriori}, a word
$\varsigma=m_{a_1}^{\eps_1}\cdots m_{a_s}^{\eps_s}$ factors as $\varsigma=b(\varsigma)\prod_{i=1}^s(1+x_i)$ with
$1+x_i=\beta_i^{-1}\,b(m_{a_i}^{\eps_i})^{-1}m_{a_i}^{\eps_i}\,\beta_i$ and
$\beta_i=b(m_{a_{i+1}}^{\eps_{i+1}}\cdots m_{a_s}^{\eps_s})\in b(\Gamma)$. Here
$b(m_a)^{-1}m_a=b_a^{-1}m_a$ and $b(m_a^{-1})^{-1}m_a^{-1}=b_am_a^{-1}$ lie in $U_P$ and normalize $S$, hence lie
in $Z$; so $n_i:=b(m_{a_i}^{\eps_i})^{-1}m_{a_i}^{\eps_i}-I\in\mathfrak z\subseteq C(S)$ has polynomial size, and
$x_i=\beta_i^{-1}n_i\beta_i=\psi_i(n_i)$ with $\psi_i=\operatorname{Ad}(\beta_i^{-1})|_{C(S)}\in\Psi$; as
$\beta_i\in L_P$ normalizes $S$, $x_i\in\mathfrak z$. Hence the $x_i$ are strictly block upper triangular with
entries in $\varrho^{-1}\Z$ of absolute value at most $2^{p_0}$, for some $\varrho$ and $p_0$ of polynomial
bit length.
Products of $m$ of them vanish, so $\varphi(\varsigma)=\prod_i(1+x_i)$ is a sum of at most $ms^{m-1}$ products of
fewer than $m$ factors $x_i$, and has bit length $O(m\log s+mp_0+m\log\varrho)$, which is polynomial since
$\log|\Psi|$ is. The same holds for $\log\varphi(\varsigma)$. Finally, the closure of the group generated by
unipotent elements $u_i\in U_P$ is $\exp\mathfrak l$, where $\mathfrak l$ is the Lie algebra generated by the
$\log u_i$: $\exp\mathfrak l$ is a closed subgroup (Baker--Campbell--Hausdorff) containing the $u_i$, and the
closure contains $\exp(\tau\log u_i)$ for all $\tau$ (A3), so its Lie algebra contains $\mathfrak l$. By (F3),
$\operatorname{Lie}(U'')$ is spanned by matrices of polynomial size.
\end{proof}

\begin{proposition}\label{prop:ugen}
Let $V$ be the smallest subspace of $\fu_P$ that contains $\log\hat x_1,\ldots,\log\hat x_k$ and
$\operatorname{Lie}(U'')$ and is stable under $\operatorname{Ad}(m_a)^{\pm1}$ ($1\le a\le k$) and under $\operatorname{ad}(Y')$ for $Y'\in\operatorname{Lie}(S)$.
Then $\fu$ is the Lie algebra generated by $V$.
\end{proposition}
\begin{proof}
Let $\mathfrak n$ be the Lie algebra generated by $V$. By (F3) it is stable under the automorphisms
$\operatorname{Ad}(m_a)^{\pm1}$ and the derivations $\operatorname{ad}(Y')$, hence under $\operatorname{Ad}(S)$ ($S$ is connected) and under
$\operatorname{Ad}(N')$, since stabilizing a subspace is a closed condition. Since $\hat x_a\in U$, $U''\subseteq U$, and
$N'\subseteq\cG$ normalizes $U$, we have $\mathfrak n\subseteq\fu$. Conversely, $\mathcal N=\exp\mathfrak n$ is a closed
connected unipotent subgroup normalized by $N'$, so $N'\mathcal N$ is a closed subgroup. It contains
$M_a=\hat x_am_a$, hence $\cG$. If $g=n'n\in U=\cG\cap U_P$ with $n'\in N'$ and $n\in\mathcal N\subseteq U_P$, then
$n'\in N'\cap U_P=U''=\exp\operatorname{Lie}(U'')\subseteq\mathcal N$, so $g\in\mathcal N$. Thus $U\subseteq\mathcal N$ and
$\fu\subseteq\mathfrak n$.
\end{proof}

\begin{proof}[Proof of Theorem~\ref{thm:detpoly}]
By Lemmas~\ref{lem:torus}, \ref{lem:xhat} and \ref{lem:Upp} ($\log\hat x_a=\sum_{i<m}(-1)^{i+1}(\hat x_a-I)^i/i$
has polynomial size), the generators of $V$ in
Proposition~\ref{prop:ugen} are spanned by matrices of polynomial size, and the operators $\operatorname{Ad}(m_a)^{\pm1}$ and
$\operatorname{ad}(Y')$ ($Y'$ in the basis $Y_1,\ldots,Y_n$ of Lemma~\ref{lem:v0}; stability under these is stability
under all of $\operatorname{ad}(\operatorname{Lie}S)$) have polynomial size. By (F2) and (F3),
$\fu$ is spanned by matrices of polynomial size, in the weight basis conjugated by $v_0$. Undoing the
conjugation by $v_0$ and the block diagonal change of basis, both of polynomial size, gives the first
claim in the adapted basis. For every $j$, $U_j=\pi_j(U)$ (Lemma~\ref{lem:radical}) and $\pi_j$ is linear, so in characteristic zero
$\fu_j=\operatorname{Lie}(U_j)=\pi_j(\fu)$; these are exactly the subspaces computed in step~3
(Proposition~\ref{prop:layer}(2)) and entering Definition~\ref{def:heights} and Lemma~\ref{lem:sizes}.
Projecting the spanning matrices of $\fu$ onto the distances $\le j$ gives spanning sets of $\fu_j$ of
polynomial size, so $\operatorname{ht}(\fu_j)$ is polynomial by (F1) for every $j$, and so is $B_f$. By
Theorem~\ref{thm:polytime}(2) the algorithm of Section~\ref{ssec:layers} then runs in deterministic polynomial
time on every invertible series.
\end{proof}

\begin{remark}\label{rem:detpoly}
(i) The algorithm is unchanged; only its analysis is new. The proof is constructive except for the spanning set
of $\operatorname{Lie}(U'')$, which uses Schreier generators of a subgroup of possibly exponential index. That set is used
only to bound heights, in the same way as in Lemma~\ref{lem:apriori}, which is the case $S=1$ of the argument
(then $Y=0$, $v_0=I$, $\hat x_a=I$, $m_a=M_a$ and $U''=U$).

(ii) Remark~\ref{rem:polytime}(ii) shows that unipotent parts of words of length $\le n$ may all be trivial in
dimension $n+2$. Lemma~\ref{lem:torus} shows that the torus behaves better: $\operatorname{Lie}(H^\circ)$ is always spanned by
the Lie algebras of the closures of the $b(M_w)$ with $|w|\le2d^2+2d$. For the affine Weyl group of type
$\tilde A_n$ in its multiplicative form of Remark~\ref{rem:polytime}(ii) (a homomorphic image of the Coxeter
group, so words of length $\le n$ still have finite order), words of length at most $n$ have
finite order, so length $n+1$ is necessary; in our exact experiments ($n\le3$) it was also sufficient.

(iii) In particular no randomness is needed in Theorem~\ref{thm:polytime}(3).
\end{remark}
